% ============================================================================
%  Bruise Segmentation -- The Complete Story
%  Compile:  tectonic docs\bruise_complete_story.tex
%  Self-contained article class. Every number traces to a shipped CSV or to
%  PROJECT_HANDBOOK.md; Appendix A lists the source for each table.
% ============================================================================
\documentclass[11pt]{article}
\usepackage[a4paper,margin=1in]{geometry}
\usepackage{booktabs}
\usepackage{array}
\usepackage{amsmath,amssymb}
\usepackage{xcolor}
\usepackage{fancyhdr}
\usepackage{enumitem}
\usepackage{longtable}
\usepackage{caption}
\usepackage{microtype}

\definecolor{secblue}{RGB}{31,73,125}
\definecolor{boxgray}{RGB}{244,246,249}
\definecolor{wingreen}{RGB}{20,110,60}
\definecolor{losered}{RGB}{150,40,40}
\definecolor{warnamber}{RGB}{150,100,10}
\definecolor{rulegray}{RGB}{200,205,212}

\usepackage[colorlinks=true,linkcolor=secblue,urlcolor=secblue,citecolor=secblue]{hyperref}
\usepackage{sectsty}
\sectionfont{\color{secblue}\large}
\subsectionfont{\color{secblue}\normalsize}
\usepackage[most]{tcolorbox}

\pagestyle{fancy}\fancyhf{}
\renewcommand{\headrulewidth}{0.4pt}
\lhead{\small Bruise Segmentation}\rhead{\small The Complete Story}\cfoot{\thepage}

\newcolumntype{n}{>{\raggedleft\arraybackslash}p{1.35cm}}
\newcolumntype{M}{>{\raggedleft\arraybackslash}p{1.90cm}}
\newcolumntype{C}{>{\raggedleft\arraybackslash}p{3.10cm}}
\newcolumntype{L}{>{\raggedright\arraybackslash}p{4.6cm}}
\newcolumntype{W}{>{\raggedright\arraybackslash}p{5.8cm}}

% Long \texttt{} tokens (run ids, file names) cannot hyphenate; \sloppy trades
% inter-word spacing for overfull boxes, which reads better than text in the
% margin. \emergencystretch gives TeX a last-resort allowance on the worst lines.
\setlength{\emergencystretch}{3em}

\newtcolorbox{keybox}[1]{colback=boxgray,colframe=secblue,boxrule=0.8pt,
  arc=2pt,left=6pt,right=6pt,top=5pt,bottom=5pt,title={\bfseries #1}}
\newtcolorbox{warnbox}[1]{colback=white,colframe=warnamber,boxrule=0.8pt,
  arc=2pt,left=6pt,right=6pt,top=5pt,bottom=5pt,title={\bfseries #1}}

\newcommand{\win}{\textcolor{wingreen}{\textbf{WIN}}}
\newcommand{\noninf}{\textcolor{secblue}{\textbf{NON-INF}}}
\newcommand{\inc}{\textcolor{warnamber}{\textbf{INCONCL}}}
\newcommand{\infr}{\textcolor{losered}{\textbf{INFERIOR}}}

\title{\vspace{-1.5cm}\color{secblue}\textbf{Bruise Segmentation in White-Light
Photography}\\[4pt]\large\color{black}The Complete Story: Dataset, Protocol,
Twelve Model Families, Fifty-Two Training Runs,\\ and What Survives Statistical
Scrutiny}
\author{Tarun Bommawar}
\date{30 July 2026}

\begin{document}
\maketitle
\thispagestyle{fancy}

\begin{keybox}{The one-paragraph version}
Binary segmentation of bruises in 1016 white-light photographs from 143
subjects. Fifty-two training runs across twelve model families. The governing
finding is that \textbf{every model sits at the annotation ceiling}: a Friedman
test across the seven headline models does not reject ($p=0.607$, Kendall
$W=0.027$), and human annotators disagree with each other by more than the
models disagree with each other. Where the models genuinely separate is
\textbf{complete-miss rate} --- three models miss zero of 185 test images, one
misses thirteen. The deployable recommendation is
\texttt{segformer\_b0\_distilled}: zero complete misses at 3.71\,M parameters,
statistically indistinguishable from its own 27\,M teacher. The single
statistically significant improvement anywhere in the study is
\texttt{lraspp\_mobilenetv3\_distilled} over its direct counterpart
($\Delta=+0.0218$ Dice, Holm-adjusted $p=0.0042$) --- and it improves the
endpoint it was \emph{not} designed to improve.
\end{keybox}

\tableofcontents
\vspace{6pt}\hrule\vspace{10pt}

% ============================================================================
\section{The problem, and why the usual metric is the wrong one}

The clinical task is not ``how well outlined is this bruise'' but \textbf{``was
the bruise found at all.''} A segmentation that traces a bruise at Dice 0.65
instead of 0.85 is clinically equivalent --- a human reviewer sees the region
and assesses it. A segmentation that returns nothing, or fires on the wrong part
of the arm, is a clinical failure of a different kind.

That asymmetry governs every reporting decision in this document. The primary
metric is \textbf{complete-miss rate}: the fraction of test images on which a
model's prediction has zero overlap with the ground truth ($\mathrm{Dice}=0$).
Mean Dice is reported alongside it, and median Dice alongside that, because the
per-image distribution is strongly left-skewed and the three tell different
stories about the same model.

\subsection{The annotation ceiling}

Three human annotators labelled overlapping subsets of the data. Their
disagreement with each other bounds what any model can be asked to achieve.

\begin{center}
\begin{tabular}{Lnn}
\toprule
\textbf{Comparison} & \textbf{mean} & \textbf{median} \\
\midrule
human: gbarimah vs majority & 0.873 & 0.926 \\
human: erik vs majority & 0.866 & 0.892 \\
\textbf{model: segformer\_b2\_teacher} & \textbf{0.769} & \textbf{0.819} \\
\textbf{model: segformer\_b0\_distilled} & \textbf{0.768} & \textbf{0.817} \\
\textbf{model: segformer\_b0\_direct} & \textbf{0.766} & \textbf{0.813} \\
human: gbarimah vs erik & 0.755 & 0.809 \\
\textbf{model: yolo\_sem\_distilled} & \textbf{0.726} & \textbf{0.801} \\
\textbf{model: yolo\_sem\_direct} & \textbf{0.702} & \textbf{0.806} \\
human: paul vs majority & 0.700 & 0.750 \\
human: paul vs gbarimah & 0.581 & 0.632 \\
human: paul vs erik & 0.581 & 0.616 \\
\bottomrule
\end{tabular}
\end{center}

The entire model field sits between \texttt{paul\_vs\_majority} (0.700) and
\texttt{gbarimah\_vs\_erik} (0.755). \textbf{A 0.005 Dice gap between two models
is not a result} --- it is inside the noise floor of the labels themselves.

\begin{keybox}{Three consequences for anything written from this study}
\begin{enumerate}[leftmargin=*,itemsep=1pt]
\item Lead with complete-miss rate, not mean Dice.
\item Never claim model X beats model Y on a Dice difference below $\sim$0.05
without a paired subject-level bootstrap whose interval excludes zero.
\item Always show the annotation ceiling alongside any model ranking.
\end{enumerate}
\end{keybox}

\subsection{Two aggregations --- never mix them}

The numbers above are the \textbf{val-selected best seed}. Elsewhere this
document quotes \textbf{3-seed mean $\pm$ std} for the same models, which is why
B2 appears as 0.769 here and $0.7651\pm0.0037$ in \S\ref{sec:stagea}. Both are
correct; every table states which it uses. The best seed is not constant across
models: seed 0 for the three SegFormers and \texttt{yolo\_sem\_distilled},
seed 2 for \texttt{yolo\_sem\_direct}.

% ============================================================================
\section{Dataset and split}

\begin{center}
\begin{tabular}{lnn}
\toprule
\textbf{Split} & \textbf{Images} & \textbf{Subjects} \\
\midrule
train & 697 & 95 \\
val & 134 & 20 \\
test & 185 & 28 \\
\midrule
total & 1016 & 143 \\
\bottomrule
\end{tabular}
\end{center}

\textbf{Subject-grouped, never image-grouped.} No subject appears in two splits
and no image appears twice. Both properties are re-asserted at bundle build time
and again in \S2 of the notebook, because a manifest is a text file and text
files get edited. A subject leaking across splits would inflate every number in
the study invisibly.

\textbf{Physical layout.} \texttt{data/train/} holds all 831 train+val images
together; a \texttt{split} column in the manifest decides which is which. Val is
a \emph{subset} of the training-side files, separated by a column and never by
directory --- duplicating val images into their own folder would create a second,
silently divergent source of truth for the 697/134 boundary.

\textbf{Metadata per image:} \texttt{stem}, \texttt{subject}, \texttt{ITA}
(individual typology angle), \texttt{ita\_group\_index\_5},
\texttt{skin\_tone\_category}, \texttt{split}, \texttt{image\_path},
\texttt{mask\_path}.

\textbf{ITA groups}, fixed order light $\rightarrow$ dark, used in every fairness
table: Light (II--III), Intermediate (III--IV), Tan (IV), Brown (V), Dark (VI).
Test-set group sizes: 39 / 38 / 24 / 29 / 55 images.

\textbf{The 640 cache.} Models train on $640\times640$ PNGs derived from
native-resolution JPEGs --- image bilinear, \textbf{mask nearest-neighbour}. That
distinction is not cosmetic: bilinear interpolation on a mask produces fractional
boundary pixels, and thresholding those back to binary erodes or dilates small
bruises, which is exactly the population the complete-miss metric is about. The
cache is derived, not shipped, and rebuilds deterministically in about five
minutes.

% ============================================================================
\section{The shared protocol}

Every model in Stages A, B, E and F trains under the \textbf{same recipe}. That
is the entire basis for comparing them: if the recipe varied, architectural
differences would be confounded with tuning effort.

\begin{center}
\begin{tabular}{lW}
\toprule
\textbf{Setting} & \textbf{Value / rationale} \\
\midrule
\texttt{img\_size} & 640 \\
\texttt{epochs} & 100, with early stopping \\
\texttt{patience} & 15 \\
\texttt{backbone\_lr} & $6\times10^{-5}$ --- pretrained encoder, conservative \\
\texttt{head\_lr} & $6\times10^{-4}$ --- random head, $10\times$, catches up \\
\texttt{betas} & (0.9, 0.999) \\
\texttt{weight\_decay} & 0.01 \\
\texttt{warmup\_fraction} & 0.01 \\
\texttt{poly\_power} & 1.0 (poly LR decay) \\
\texttt{gradient\_clip} & 1.0 \\
\texttt{amp} & True \\
\texttt{seeds} & (0, 1, 2) \\
\bottomrule
\end{tabular}
\end{center}

\textbf{Why the $10\times$ head LR.} The backbone is pretrained and already has
good features; a conservative rate preserves them. The head is randomly
initialised and must catch up. This is the SegFormer recipe (Xie et al.\ 2021),
applied to \emph{every} architecture here --- including YOLO, not because YOLO's
paper says so, but because holding the recipe fixed across architectures is the
whole point.

\textbf{What the LR split actually splits.} \texttt{build\_param\_groups}
assigns groups by \texttt{id(p)} against \texttt{model.backbone}, not by name ---
the architectures name their parameters completely differently and a name-prefix
rule would put every YOLO parameter in the wrong group. One consequence worth
knowing: Fast-SCNN's \texttt{.backbone} is only \texttt{learning\_to\_downsample},
three conv blocks. Its global feature extractor, feature fusion and classifier
all land in the head group. Nothing in Fast-SCNN is pretrained, so the split is
provenance-neutral --- but those three convs genuinely do learn $10\times$ slower
than everything around them. It applies identically to \texttt{fastscnn} and
\texttt{fastscnn\_distilled}, so the Stage F contrast is unaffected; only
Fast-SCNN's absolute number is.

\textbf{Loss:} Dice + BCE (\texttt{SupervisedLoss}), or \texttt{DistillLoss} when
a teacher is present. \texttt{aux\_weight}~=~0.4 for SegFormer; 0.0 for baselines
and mobile models, which have no auxiliary head.

\textbf{Model selection:} best validation \textbf{AP} (threshold-free), not best
Dice at some threshold. Selecting on a thresholded metric would entangle model
choice with threshold choice.

\subsection{Threshold selection --- the operating point}

After training, the logit cut is swept over 481 values on the 134 validation
images, then applied \emph{once} to test.

The cut is \textbf{not} the argmax. These sweeps are extraordinarily flat: B2's
validation Dice moved by 0.009 across thresholds from 0.154 to 0.959. That is not
a peak, it is noise on a plateau, and taking the argmax fits the validation set's
sampling error. \texttt{select\_cut} therefore treats every cut within
\textbf{one standard error} of the peak as statistically tied, and breaks the tie
by \textbf{lowest complete-miss rate}. Cuts in the band are Dice-equivalent but
they are \emph{not} miss-equivalent.

A checkpoint without its \texttt{operating\_point.json} is unusable --- its test
score is undefined. The registry enforces this.

\subsection{The normalisation contract}

The dataloader emits \textbf{raw $[0,1]$ pixels} and each model applies its own
scale internally. This is not a style preference. SegFormer wants ImageNet
normalisation; YOLO wants plain $/255$, and Ultralytics' BatchNorms carry frozen
running statistics for the $/255$ distribution. Feeding YOLO ImageNet-normalised
pixels makes it under-fire by $4\times$ and \textbf{caps it at Dice 0.479 with no
threshold able to recover it}. Pixel scale belongs to the model, not the loader.

% ============================================================================
\section{Stage A --- the five headline models}
\label{sec:stagea}

\begin{center}
\begin{tabular}{LnW}
\toprule
\textbf{Model} & \textbf{Params} & \textbf{Description} \\
\midrule
\texttt{segformer\_b2\_teacher} & 27.35\,M & SegFormer MiT-B2, the teacher \\
\texttt{segformer\_b0\_direct} & 3.71\,M & SegFormer MiT-B0, supervised only \\
\texttt{segformer\_b0\_distilled} & 3.71\,M & B0 distilled from the same-seed B2 \\
\texttt{yolo\_sem\_direct} & 1.63\,M & YOLO26n-seg, native Ultralytics \\
\texttt{yolo\_sem\_distilled} & 1.63\,M & YOLO26n-seg, offline pseudo-mask KD \\
\bottomrule
\end{tabular}
\end{center}

Three seeds each: 15 runs, all shipped.

\subsection{Results (3-seed mean $\pm$ std, 185 test images)}

\begin{center}
\begin{tabular}{LMnM}
\toprule
\textbf{Variant} & \textbf{mean Dice} & \textbf{median} & \textbf{miss \%} \\
\midrule
segformer\_b2\_teacher & $0.7651\pm0.0037$ & 0.8113 & $0.00\pm0.00$ \\
segformer\_b0\_distilled & $0.7639\pm0.0047$ & 0.8114 & $0.18\pm0.31$ \\
segformer\_b0\_direct & $0.7600\pm0.0055$ & 0.8113 & $0.54\pm0.00$ \\
yolo\_sem\_direct (native) & $0.6911\pm0.0197$ & 0.8021 & $7.57\pm1.08$ \\
yolo\_sem\_distilled (native) & $0.6676\pm0.0648$ & 0.7658 & $7.57\pm4.80$ \\
\bottomrule
\end{tabular}
\end{center}

\textbf{Distillation works, barely.} B0-distilled gains $+0.004$ Dice over
B0-direct --- statistically indistinguishable --- but drops complete misses from
0.54\,\% to 0.18\,\%. \textbf{The miss-rate improvement is the real effect; the
Dice difference is not.} This is the template for every distillation claim in
this study.

\subsection{SegFormer distillation details}

\begin{itemize}[leftmargin=*,itemsep=1pt]
\item \textbf{Teacher: the same-seed B2, not the best B2.} Pairing seed $k$'s
student with seed $k$'s teacher means the reported spread includes the teacher's
own variance, which is part of the pipeline being measured. Using one strong
teacher for all three students would make the spread artificially narrow.
\item Teacher temperature is \textbf{calibrated} on validation after training
(\texttt{calibrate\_temperature}), written to \texttt{calibration.json}.
Measured $T\approx1.69$--$1.76$.
\item \texttt{segformer\_alpha} = 0.6, the KD mix.
\end{itemize}

\subsection{YOLO: the two-path story}

YOLO is evaluated two ways, and only one is a reporting path.

\textbf{\texttt{native\_argmax}} --- Ultralytics' own semantic head, argmax over
classes. Parameter-free: there is no threshold to fit, so nothing to overfit.
\textbf{This is the sole reporting path.}

\textbf{\texttt{custom255}} --- a custom probability path with a swept threshold.
Retained in the bundle for provenance only. It must not appear in charts.

YOLO trains \textbf{natively} (mosaic, EMA, letterbox, its own LR schedule)
rather than through the shared custom loop. Distilled YOLO uses \textbf{offline
pseudo-masks}: the same-seed B2 teacher's native-resolution probability fused
with ground truth as
$\mathrm{class} = (\alpha\,\mathrm{GT} + (1-\alpha)\,p_{\mathrm{teacher}} \ge 0.5)$
with \texttt{yolo\_alpha} = 0.4. Alpha below 0.5 keeps the fusion non-degenerate.
Pseudo-masks are built at native resolution because Ultralytics letterboxes
internally --- pre-resizing would apply the resize twice.

\subsection{Speed (full A100, 640 tensor $\rightarrow$ mask on GPU)}

\begin{center}
\begin{tabular}{Lnnnn}
\toprule
\textbf{Model} & \textbf{ms} & \textbf{FPS} & \textbf{params M} & \textbf{act.\ MB} \\
\midrule
segformer\_b2\_teacher & 33.67 & 29.7 & 27.35 & 1778.7 \\
segformer\_b0\_direct & 16.68 & 59.9 & 3.71 & 1204.4 \\
segformer\_b0\_distilled & 16.64 & 60.1 & 3.71 & 1204.4 \\
yolo\_sem\_direct & 8.18 & 122.2 & 1.63 & 967.5 \\
yolo\_sem\_distilled & 8.22 & 121.7 & 1.63 & 967.5 \\
\bottomrule
\end{tabular}
\end{center}

Disk read, JPEG decode, resize and host/GPU copies are deliberately \textbf{not}
timed --- they are identical for every model and dominated by I/O, so including
them would compress real architectural differences into measurement noise.
\texttt{cuda.synchronize()} brackets every call; without it the benchmark
measures how long it takes to \emph{queue} the work and reports every model as
equally fast.

\subsection{Fairness (Stage A, ITA groups)}

\begin{center}
\begin{tabular}{Lnlllr}
\toprule
\textbf{Model} & \textbf{KW $p$} & \textbf{sig.} & \textbf{gap} & \textbf{worst} & \textbf{miss gap} \\
\midrule
segformer\_b2\_teacher & 0.0105 & \textbf{yes} & 0.112 & Tan & 0.000 \\
segformer\_b0\_direct & 0.639 & no & 0.052 & Dark & 0.000 \\
segformer\_b0\_distilled & 0.470 & no & 0.050 & Dark & 0.000 \\
yolo\_sem\_direct & 0.230 & no & 0.104 & Light & 0.154 \\
yolo\_sem\_distilled & 0.151 & no & 0.090 & Light & 0.051 \\
\bottomrule
\end{tabular}
\end{center}

\textbf{The counter-intuitive result:} where YOLO shows a fairness gap, the
\emph{worst} group is Light (II--III), not Dark. Light-skin under-detection, the
opposite of the expected direction. See \S\ref{sec:sizeconfound} for why this is
probably a size confound and must not be reported as a skin-tone effect without
conditioning on bruise size.

% ============================================================================
\section{Stage B --- direct baselines}

Same recipe, ImageNet ResNet-50 encoders, fresh 1-class heads. Three seeds each,
six runs, all shipped.

\begin{center}
\begin{tabular}{LMnM}
\toprule
\textbf{Variant} & \textbf{mean Dice} & \textbf{median} & \textbf{miss \%} \\
\midrule
unet\_r50 & $0.7526\pm0.0154$ & 0.826 & $3.24\pm0.54$ \\
deeplabv3plus\_r50 & $0.7453\pm0.0182$ & 0.814 & $1.80\pm0.62$ \\
\bottomrule
\end{tabular}
\end{center}

\textbf{nnU-Net was never run on the canonical split.} No weights, no results. It
is registered as an explicit gap rather than omitted, because ``the baseline we
did not run'' is exactly the fact a reader needs. Enabling it needs
\texttt{nnunetv2} and roughly 8 GPU-hours.

\textbf{Fairness:} neither baseline shows a significant ITA effect
(U-Net $p=0.222$, DeepLab $p=0.108$).

\begin{keybox}{U-Net is the study's clearest metric trap}
U-Net has the \textbf{highest median Dice of any model in the study} (0.8329 at
its best seed) and 7 complete misses out of 185. SegFormer-B2 has a lower median
(0.8192) and \textbf{zero} misses. Rank on median and U-Net wins; rank on the
clinical metric and it is in the bottom tier. Both numbers are correct. Only one
of them answers the question.
\end{keybox}

% ============================================================================
\section{Stage C --- SegFormer-B5 distillation}

The most elaborate stage. A B5 teacher was trained, then a grid of KD strategies
distilled B5 (and B2+B5 ensembles) into B0 students, each scored against the
\textbf{B2$\rightarrow$B0 reference} from Stage A.

\subsection{The B5 teacher}

Trained at three seeds; the validation-winning seed was promoted to
\texttt{teachers/}.

\begin{center}
\begin{tabular}{lnnnn}
\toprule
\textbf{seed} & \textbf{test mean} & \textbf{median} & \textbf{threshold} & \textbf{misses} \\
\midrule
42 & 0.7527 & 0.8124 & 0.35 & 0 \\
\textbf{123 (selected)} & \textbf{0.7727} & \textbf{0.8273} & \textbf{0.50} & \textbf{0} \\
2026 & 0.7676 & 0.8298 & 0.40 & 0 \\
\bottomrule
\end{tabular}
\end{center}

Selected on \textbf{validation} (\texttt{val\_mean\_dice}: 0.7787 / 0.7881 /
0.7861), never on test. \texttt{segformer\_b5\_teacher} \emph{is}
\texttt{segformer\_b5\_\_seed123} --- the same weights under two names.
\texttt{report.load\_stage\_c} detects the identical per-image vectors and
collapses them, because listing both would show one result twice and invite
reading it as two agreeing runs.

\subsection{The multi-teacher gate --- \texttt{val\_oracle}}

\emph{Before} running any multi-teacher arm, an oracle test asked whether B2 and
B5 are complementary enough to be worth combining. Run on the 134 validation
images:

\begin{center}
\begin{tabular}{lC}
\toprule
Spearman $\rho$(B2, B5) per-image Dice & 0.819 \\
B5 beats B2 by $>0.10$ & 18 images \\
B2 beats B5 by $>0.10$ & 12 images \\
B2 mean Dice / B5 mean Dice & 0.7717 / 0.7881 \\
\textbf{oracle mean Dice} & \textbf{0.8140} \\
oracle gain over best single & $+0.0258$ \\
95\,\% CI of gain & $[0.0165, 0.0362]$ \\
$P(\text{gain} > 0)$ & 1.00 \\
\texttt{rel\_b} for adaptive gate & 0.6343 \\
\textbf{GATE\_run\_adaptive} & \textbf{true} \\
\bottomrule
\end{tabular}
\end{center}

The two teachers correlate at $\rho=0.82$ but disagree substantially on 30 of 134
images. An oracle picking the better teacher per image would gain $+0.026$ Dice
over the best single teacher. \textbf{The gate opened}, and $\mathrm{rel\_b}=0.634$
became the adaptive arms' mixing weight. This is the right shape for an
experiment: a cheap, decisive pre-test that determines whether an expensive grid
is worth running at all.

\subsection{Alpha search}

The KD mix was searched on validation with 15-epoch short runs, 5 grid points
(Optuna if installed, grid otherwise). Ten short runs total; they produce no
model and are not counted among the 52 real runs.

\begin{center}
\begin{tabular}{lnnW}
\toprule
\textbf{Tag} & \textbf{best $\alpha$} & \textbf{val Dice} & \textbf{trials ($\alpha\to$ val Dice)} \\
\midrule
\texttt{single\_b5\_response} & 0.625 & 0.7019 & 0.300$\to$0.7017, 0.462$\to$0.6951, 0.625$\to$0.7019, 0.787$\to$0.6975, 0.950$\to$0.7010 \\
\texttt{ensemble\_uniform} & 0.950 & 0.7089 & 0.300$\to$0.7016, 0.462$\to$0.7067, 0.625$\to$0.7004, 0.787$\to$0.6952, 0.950$\to$0.7089 \\
\bottomrule
\end{tabular}
\end{center}

Note how flat both are --- 0.007 across the whole range for single-teacher. This
is the same plateau phenomenon as the threshold sweep. \textbf{Alpha is not a
sensitive knob here}, and the study should not describe it as ``tuned'' in any
meaningful sense.

\subsection{The arms}

All students are SegFormer-B0, seed 42, on the 697/134 split.

\begin{center}
\begin{tabular}{LnnW}
\toprule
\textbf{Arm} & \textbf{mean} & \textbf{$\alpha$} & \textbf{What it does} \\
\midrule
\texttt{p3\_adaptive} & \textbf{0.7748} & 0.95 & B2+B5 adaptive ensemble, per-image gating at rel\_b=0.634 \\
\texttt{p2\_cwd\_b5\_to\_b0} & 0.7736 & 0.625 & Channel-wise distillation from B5 \\
\texttt{expA\_b5\_to\_b0\_response} & 0.7683 & 0.625 & Plain response KD from B5 \\
\texttt{p3\_adaptive\_boundary} & 0.7668 & 0.95 & Adaptive + boundary loss ($\lambda=4.0$) \\
\texttt{p2\_bpkd\_b5\_to\_b0} & 0.7651 & 0.625 & Boundary-privileged KD \\
\texttt{p2\_ensemble\_uniform} & 0.7615 & 0.95 & B2+B5 uniform ensemble \\
\texttt{p3\_adaptive\_group} & 0.7586 & 0.95 & Adaptive + ITA-group weighting ($\lambda=0.5$) \\
\texttt{p3\_adaptive\_full} & 0.7515 & 0.95 & Adaptive + all auxiliary terms \\
\texttt{x\_angular\_b5\_to\_b0} & 0.7418 & 0.625 & Angular/relational KD ($\lambda=1.0$) \\
\texttt{p3\_adaptive\_hard} & 0.7414 & 0.95 & Adaptive + hard-example focus ($\gamma=2.0$) \\
\texttt{x\_dkd\_b5\_to\_b0} & --- & 0.625 & \textbf{Never executed} (config only) \\
\bottomrule
\end{tabular}
\end{center}

\subsection{How arms are scored, and the honest summary}

Each arm is compared to the B2$\rightarrow$B0 reference (0.7680) with a paired
subject-level bootstrap and a non-inferiority margin of one Dice point. At 28
subjects, \textbf{every arm came back NON-INFERIOR}. \texttt{p3\_adaptive} leads
with $\Delta=+0.0068$, CI $[-0.0025, +0.0152]$, $P(\text{better})=0.93$ ---
suggestive, not significant.

\begin{keybox}{The honest summary of Stage C}
None of the ten KD strategies produced a statistically distinguishable
improvement over plain B2$\rightarrow$B0 distillation. That is a real result and
should be reported as one, not buried. See \S\ref{sec:inconclusive} for a
correction to how these verdicts were assigned.
\end{keybox}

% ============================================================================
\section{Stage E --- mobile baselines}

Four deployment-scale architectures, trained directly under the same recipe.

\begin{center}
\begin{tabular}{LnW}
\toprule
\textbf{Model} & \textbf{Params} & \textbf{Init / source} \\
\midrule
PP-MobileSeg-Tiny & 1.45\,M & StrideFormer-Tiny backbone, ImageNet (OpenMMLab, auto) \\
TopFormer-Tiny & 1.37\,M & ImageNet backbone, 66.2\,\% top-1 (Google Drive, manual) \\
LR-ASPP MobileNetV3 & 3.22\,M & MobileNetV3-Large, ImageNet (torchvision, auto) \\
Fast-SCNN & 1.14\,M & \textbf{none --- scratch by design} \\
\bottomrule
\end{tabular}
\end{center}

\textbf{Fast-SCNN's lack of weights is not a gap.} At 1.1\,M parameters the paper
trains from scratch and no official checkpoint was ever released; random
initialisation is the faithful reproduction. But it does mean Fast-SCNN is not
initialised like the others, and that must be stated wherever it is compared.

\subsection{Vendoring, and why}

PP-MobileSeg and TopFormer only have reference implementations inside the
OpenMMLab stack. \texttt{mmcv} ships compiled CUDA extensions pinned to narrow
torch versions and routinely fails to build; dragging it in for two 1.4\,M
parameter models would be the heaviest dependency in the project.

But \textbf{hand-rewriting them is worse}. The published checkpoints are state
dicts keyed by module path --- any deviation in naming makes
\texttt{load\_state\_dict} fail or, worse, half-succeed and leave most weights
random while reporting success.

So the reference files are \textbf{vendored verbatim} into
\texttt{bruisekit/vendor/}, with only their import block rewritten to
\texttt{bruisekit/mmcv\_shim.py}. The shim reproduces \texttt{ConvModule}
(including the attribute names \texttt{conv}, \texttt{bn}, \texttt{activate},
which \emph{are} the state-dict keys), \texttt{build\_norm\_layer},
\texttt{build\_activation\_layer}, \texttt{DropPath}, and inert stubs for the
registry plumbing.

\begin{warnbox}{One substitution that would have been silent}
StrideFormer's SE blocks ask for
\texttt{dict(type='Hardsigmoid', slope=0.2, offset=0.5)}.
\texttt{nn.Hardsigmoid} fixes the slope at $1/6$. Substituting the builtin leaves
every state-dict key matching --- this layer has no parameters --- while quietly
changing the gate on every SE block in the backbone. The shim implements
\texttt{HSigmoid} with honoured slope and offset.
\end{warnbox}

Verification that the vendoring worked:

\begin{center}
\begin{tabular}{Lnnnl}
\toprule
\textbf{arch} & \textbf{loaded} & \textbf{in ckpt} & \textbf{unexp.} & \textbf{verdict} \\
\midrule
ppmobileseg\_tiny & 488 & 488 & 0 & EXACT MATCH \\
lraspp\_mobilenetv3 & 308 & 308 & 0 & EXACT MATCH \\
\bottomrule
\end{tabular}
\end{center}

$\text{unexpected}=0$ with only the decode head missing is proof the architecture
is the published one. This matters more than parameter count: PP-MobileSeg
measures 1.454\,M against a published 1.61\,M, but that figure is for a 150-class
ADE20K head, and the checkpoint match settles it.

\textbf{Backbone-only, always.} Segmentation-pretrained checkpoints exist for
both PP-MobileSeg (ADE20K) and LR-ASPP (COCO-with-VOC-labels) and are
deliberately \textbf{not} used --- they would give those models segmentation
pretraining that no other baseline in the study has.

\subsection{Results (val-selected seed, 185 test images)}

\begin{center}
\begin{tabular}{LnnnnM}
\toprule
\textbf{Model} & \textbf{mean} & \textbf{median} & \textbf{prec} & \textbf{rec} & \textbf{misses} \\
\midrule
lraspp\_mobilenetv3\_distilled & 0.7200 & 0.7836 & 0.8398 & 0.6833 & 3 (1.62\,\%) \\
lraspp\_mobilenetv3 & 0.6982 & 0.7607 & 0.8322 & 0.6548 & 2 (1.08\,\%) \\
topformer\_tiny & 0.6918 & 0.7418 & 0.7421 & 0.7064 & \textbf{1 (0.54\,\%)} \\
ppmobileseg\_tiny & 0.6568 & 0.7158 & 0.7605 & 0.6583 & 4 (2.16\,\%) \\
fastscnn & 0.6053 & 0.6866 & 0.7601 & 0.5754 & 13 (7.03\,\%) \\
\bottomrule
\end{tabular}
\end{center}

\subsection{Reading Stage E}

\textbf{The ordering is stable, the gaps mostly are not.} LR-ASPP $>$ TopFormer
$>$ PP-MobileSeg $>$ Fast-SCNN holds on all three seeds. But LR-ASPP beats
TopFormer by only 0.006 at the val-selected seed and the paired bootstrap returns
INCONCLUSIVE (\S\ref{sec:stageg}). The LR-ASPP $\rightarrow$ Fast-SCNN gap
($+0.093$) \emph{is} large enough, and it is a \win.

\textbf{Precision and recall separate the field more cleanly than Dice.} LR-ASPP
runs at precision $\approx$ 0.83 with recall $\approx$ 0.65; TopFormer at
precision $\approx$ 0.74 with recall $\approx$ 0.71. They land at similar Dice by
different routes --- LR-ASPP is conservative and accurate where it fires,
TopFormer fires more widely. For a clinical find-the-bruise task that difference
matters more than the 0.006 Dice gap does.

\textbf{TopFormer has the best miss rate in the mobile field} --- 1 of 185, at
1.37\,M parameters. If complete-miss rate is the primary metric, and this document
argues throughout that it is, TopFormer is arguably the better mobile model
despite ranking below both LR-ASPP variants on Dice.

\textbf{PP-MobileSeg is the \emph{slowest} of the four} despite being nearly the
smallest --- its strided SEA attention is memory-bound rather than FLOP-bound.

\textbf{Fast-SCNN is fastest and least accurate}, consistent with being the only
model without pretrained initialisation. It also has the widest seed spread and
the most extreme thresholds ($-2.0$ to $-2.3$, against $\approx-0.7$ for SegFormer
and LR-ASPP): it has to push the operating point a long way negative to keep the
miss rate down, which is what a scratch-initialised model with weak features looks
like at the threshold stage.

\subsection{Speed (A100 MIG 3g.40gb slice)}

\begin{center}
\begin{tabular}{Lnnnn}
\toprule
\textbf{Model} & \textbf{ms} & \textbf{FPS} & \textbf{params M} & \textbf{act.\ MB} \\
\midrule
fastscnn & 3.56 & 281.3 & 1.135 & 965.1 \\
lraspp\_mobilenetv3 & 5.01 & 199.8 & 3.218 & 1009.3 \\
topformer\_tiny & 6.19 & 161.6 & 1.370 & 1000.9 \\
ppmobileseg\_tiny & 10.67 & 93.8 & 1.454 & 999.2 \\
\bottomrule
\end{tabular}
\end{center}

\begin{warnbox}{These are NOT comparable to the Stage A speed table}
Stage A was measured on a full A100; Stage E on a MIG 3g.40gb slice, roughly
$3/7$ of a card. Any latency claim that spans the two tables is invalid. Re-time
both on one device before publishing one.
\end{warnbox}

% ============================================================================
\section{Stage F --- DeepLabV3+ $\rightarrow$ mobile-student distillation}

\begin{sloppypar}
Two further families trained with a frozen \textbf{DeepLabV3+/ResNet-50} teacher:
\texttt{fastscnn\_distilled} and \texttt{lraspp\_mobilenetv3\_distilled}. This is
the study's second distillation axis and the first that crosses architecture
families.
\end{sloppypar}

Each arm's contrast is distilled-vs-direct on the same architecture, both at
seeds 0/1/2 under the identical recipe. Nothing else changes between the two
variants, so any difference is attributable to the teacher signal. Across arms,
teacher and recipe are also identical, so the Fast-SCNN vs LR-ASPP pair isolates
one further variable: \textbf{what the student's initialisation (scratch vs
ImageNet) does to KD}.

\subsection{Why DeepLabV3+ is the teacher}

Three reasons, in the order they matter.

\textbf{Licence.} The SegFormer MiT weights are NVIDIA's, \texttt{license: other},
non-commercial. DeepLabV3+/R50 reaches the study through
\texttt{segmentation\_models\_pytorch} (MIT) on a torchvision ResNet-50 (BSD-3),
so a student distilled from it carries no non-commercial link. This does not by
itself make the result commercialisable --- the ImageNet provenance of the
encoder and, far more importantly, the consent scope of the photographs are
separate questions --- but it removes the one unambiguously non-commercial
licence.

\textbf{Complete misses, not Dice.} DeepLabV3+ and U-Net are indistinguishable on
Dice (0.7584 vs 0.7570, with U-Net's median actually higher) --- inside the noise
floor. They are \emph{not} indistinguishable on complete misses, and for a teacher
that is the metric that decides. A teacher that misses a bruise does not hand the
student a weak soft target; it hands it a confidently-empty one on exactly the
image the clinical metric cares about. Paired over all 185 test images:

\begin{center}
\begin{tabular}{ln}
\toprule
images DeepLab misses that U-Net finds & 0 \\
images U-Net misses that DeepLab finds & 2 \\
both miss & 5 \\
\bottomrule
\end{tabular}
\end{center}

Strict containment, and the ordering holds at every seed (DeepLab 5/5/2 misses,
U-Net 9/7/7).

\textbf{Structure.} DeepLabV3+'s ASPP-then-fuse-one-low-level-skip decoder is the
same shape as Fast-SCNN's PPM-then-feature-fusion, so a later feature-level arm
has an obvious place to attach. U-Net's four symmetric skips have no counterpart
in the student.

\subsection{Why the teacher must be calibrated first}

\texttt{load\_teacher} divides the teacher's logits by a fitted temperature. The
Stage B baselines were trained as endpoints, not as teachers, so they ship no
\texttt{calibration.json} and the shim would raise.
\texttt{ensure\_calibration} fits it with the same
\texttt{engine.calibrate\_temperature} (Guo et al.\ 2017, L-BFGS on $\log T$ over
the 134 validation images) that the B2 teacher used, so both teachers are
prepared identically.

\textbf{Distilling from an uncalibrated teacher is the failure mode this guards
against}: a near-binary soft label is the hard label with extra steps, and the arm
would silently measure nothing. The shim raises rather than falling back.

Calibration is written to
\texttt{WORK\_DIR/teacher\_calibration/}, never next to the checkpoint:
\texttt{checkpoints/baselines/} is the verified artefact that produced the
published Stage B table, and adding a file to it would make the bundle differ from
the one those numbers came from.

\subsection{Settings that matter}

\texttt{efficient\_alpha = 0.6} --- the KD mix, \textbf{not}
\texttt{segformer\_alpha}. Both are currently 0.6, so reading the wrong one would
look right and silently ignore the other the moment someone tuned one of them.

\texttt{teacher\_seed\_mode = "same"} --- student seed $k$ distils from teacher
seed $k$, Stage A's convention, so the student's reported spread includes the
teacher's own variance.

\subsection{The LR-ASPP arm --- results}
\label{sec:lraspparm}

\begin{center}
\begin{tabular}{lnnnnnM}
\toprule
\textbf{seed} & \textbf{cut} & \textbf{mean} & \textbf{median} & \textbf{IoU} & \textbf{prec} & \textbf{misses} \\
\midrule
0 & $-1.275$ & 0.7406 & 0.7977 & 0.6139 & 0.8050 & 1 \\
1 & $-0.900$ & 0.7182 & 0.7782 & 0.5899 & 0.8094 & 2 \\
2 & $-1.775$ & 0.7200 & 0.7836 & 0.5912 & 0.8398 & 2 \\
\bottomrule
\end{tabular}
\end{center}

\begin{center}
\begin{tabular}{LMMM}
\toprule
& \textbf{direct} & \textbf{distilled} & \textbf{$\Delta$} \\
\midrule
3-seed mean Dice & $0.7093\pm0.0127$ & $0.7263\pm0.0125$ & $+0.0170$ \\
3-seed mean median & 0.7723 & 0.7865 & $+0.0142$ \\
val-selected seed (both s2) & 0.6982 & 0.7200 & $+0.0218$ \\
complete misses (of 555) & 3 & 5 & $+2$ \\
val Dice, seed 2 & 0.7277 & 0.7348 & $+0.0071$ \\
\bottomrule
\end{tabular}
\end{center}

Paired against direct at the same seed: $+0.0175$ / $+0.0116$ / $+0.0218$.
\textbf{The sign is the same at every seed.}

\begin{keybox}{The Stage G verdict: \win}
$\Delta = +0.0218$, 95\,\% CI $[+0.0081, +0.0372]$ --- the interval excludes zero
--- two-sided $p=0.0014$, \textbf{Holm-adjusted $p=0.0042$}. Seed consistency
3/3 positive, with seeds 0 and 2 individually excluding zero. The
\texttt{mobile\_field} omnibus rejects first ($\chi^2=16.26$, $p=0.0027$), so the
pairwise test is licensed rather than fished for. This is the only sub-0.05 Dice
difference in the entire study that clears the bar set in \S1, and that is
\emph{not} a licence to quote the others.
\end{keybox}

\textbf{What survives beyond the point estimate:}

\begin{itemize}[leftmargin=*,itemsep=2pt]
\item \textbf{Precision and recall both rose} (seed 2: $0.8322\to0.8398$ and
$0.6548\to0.6833$). A pure threshold shift trades one against the other. Both
moving up means the pixel \emph{ranking} improved --- a claim about the model, not
about the operating point.
\item \textbf{The student closed roughly half the gap to its teacher.}
DeepLabV3+/R50 sits at 0.7453; direct LR-ASPP 0.7093 (gap 0.036), distilled
0.7263 (gap 0.019), at 3.22\,M parameters against the teacher's $\sim$26.7\,M. The
student-vs-teacher contrast is $\Delta=-0.0384$, CI $[-0.0912, +0.0027]$ ---
\inc, i.e.\ the $8\times$ smaller student is \emph{not} shown to be worse.
\end{itemize}

\textbf{The pre-registered claim did not land, and the stated risk did.} The arm
was specified as a miss-rate and fairness claim, \emph{not} Dice --- because
LR-ASPP direct already sits at the annotation ceiling. Dice is the only thing that
moved in the right direction. On misses: $\Delta = +0.54$\,pp,
$P(\text{fewer misses}) = 0.19$, no containment either way (2 images the distilled
model misses that the direct one finds, 1 the other way). The DeepLab teacher
misses more than this student does (12 of 555 vs 3), so the student inheriting some
of the teacher's miss behaviour along with its Dice is the predicted failure mode,
written down before the arm ran.

\textbf{The cuts moved a long way negative} --- $-1.275$ / $-0.900$ / $-1.775$
against direct's $-0.625$ / $-0.825$. The expected reading is that a
temperature-softened teacher flattens the student's logit distribution. The
concerning reading is that the spread across seeds (0.875 of logit, against
direct's 0.2) means there is no single operating point that transfers across
seeds, which is a deployment problem even though it is not a Dice problem.

\subsection{The Fast-SCNN arm --- and why the pair matters}

The Fast-SCNN arm's previous-lineage numbers are null-to-negative:
$\Delta$ mean Dice $-0.005$ (CI $[-0.046, +0.032]$, $P(\text{better})=0.44$),
complete misses $7\to13$, and the skin-tone fairness gap widened from 0.136 to
0.199 with Kruskal-Wallis $p=0.0017$, worst group Dark (VI).

\begin{keybox}{The two-arm finding}
Fast-SCNN (scratch student): $\Delta -0.005$, misses $7\to13$, fairness gap
widens. LR-ASPP (ImageNet student): $\Delta +0.022$ significant, misses roughly
unchanged. The two arms differ only in the student, so \textbf{student
initialisation modulates the direction of cross-architecture knowledge
distillation} --- it hurts a scratch student and helps a pretrained one. Neither
arm alone could have said this. It is the most interesting claim available from
Stage F, and it is the one thing still blocked (see \S\ref{sec:next}).
\end{keybox}

\begin{warnbox}{\texttt{fastscnn\_distilled} has no current-lineage result}
All three seeds appear in \texttt{gaps.csv}. The
\texttt{per\_image\_fastscnn\_distilled.csv} on disk is from the superseded
2026-07-29 sweep. Stage G correctly \emph{skipped and named} both of its
contrasts rather than scoring stale data. Until it is re-run, the two-arm claim
above rests on cross-lineage comparison and must not be published as-is.
\end{warnbox}

% ============================================================================
\section{Stage D --- analysis methodology}

Everything in Stage D is a function of \textbf{per-image Dice}, not of model
weights. That is why the whole analysis reproduces on a laptop with no GPU and no
\texttt{torch}.

\subsection{The common schema}

Five different producers wrote per-image CSVs with different columns.
\texttt{report.normalize()} reduces all of them to one schema by
\textbf{recomputing} the derived columns from the seven-column core rather than
trusting whichever ones happen to be present:

\begin{align*}
\mathrm{tp} &= \mathrm{dice}\times(\mathrm{pred\_pos} + \mathrm{gt\_pos})/2
  \quad\text{(definition of Dice)}\\
\mathrm{fp} &= \mathrm{pred\_pos} - \mathrm{tp}, \qquad
\mathrm{fn} = \mathrm{gt\_pos} - \mathrm{tp}\\
\mathrm{complete\_miss} &= (\mathrm{dice} = 0), \qquad
\mathrm{pred\_gt\_ratio} = \mathrm{pred\_pos}/\mathrm{gt\_pos}
\end{align*}

Verified exact against the shipped outputs to $4\times10^{-12}$. Recomputing is
safer than reading: if one producer's \texttt{complete\_miss} ever disagreed with
its own \texttt{dice}, taking the stored column would propagate that bug into a
cross-model comparison. Subject and ITA columns are \textbf{always} dropped and
re-joined from the manifest, so there is one source of truth for who each image
belongs to.

\subsection{Subject-level cluster bootstrap}

The 185 test images come from \textbf{28 subjects}, and images of the same bruise
are strongly correlated. Resampling images would treat 185 correlated observations
as independent and produce intervals roughly $\sqrt{185/28}\approx2.6\times$ too
narrow.

Every CI and contrast resamples \textbf{subjects} with replacement, taking all of
a chosen subject's images. Contrasts are \textbf{paired} --- both models saw the
same 185 images, and the same resampled subject list is applied to both on each
draw; resampling independently would discard the pairing and hide real small
effects.

$P(A\text{ better})$ is reported instead of a $p$-value in Stage D: at 28 subjects
it is the honest quantity and it avoids a significant/not-significant dichotomy.

\subsection{The size $\leftrightarrow$ fairness confound}
\label{sec:sizeconfound}

Bruise size is the strongest single predictor of whether a model finds a bruise at
all, \textbf{and} it is not evenly distributed across ITA groups in this dataset.
Any fairness claim that does not condition on size is measuring both at once.

This is why the size figures sit immediately after the fairness figures rather
than in an appendix, and why the apparent Light-skin under-detection in YOLO must
not be reported as a skin-tone effect without the size-stratified analysis.

% ============================================================================
\section{Stage G --- final significance}
\label{sec:stageg}

Stage D describes; Stage G decides. They are separate sections backed by separate
modules because a confirmatory analysis only means something if its comparison
list was fixed \emph{before} the numbers were seen.

\subsection{Why not an all-pairs tournament}

Twenty scored models is 190 ordered pairs. At 28 subjects, with every model inside
the annotation-ceiling band, an all-pairs sweep at $\alpha=0.05$ is expected to
return about ten ``significant'' differences from noise alone --- and would be read
as a ranking, which is precisely what \S1 forbids.

\texttt{CONTRAST\_FAMILY} is twelve comparisons, each attached to a question
someone actually asked. Four are \textbf{confirmatory} (specified as the point of
an experiment before it ran) and Holm--Bonferroni is applied within that set of
four. The other eight are \textbf{exploratory}, reported uncorrected and labelled.

\subsection{The omnibus gates the pairwise table}

A Friedman test across models on \textbf{subject-mean} Dice --- 28 rows, not 185.
Blocking on images would reuse the same information $\sim6.6\times$ per subject.
One test over all twenty models would carry 19 degrees of freedom on a pool where
most models are B0 students of one teacher differing by a loss term, so the
omnibus runs on four pre-specified fields.

\begin{center}
\begin{tabular}{lnnnnl}
\toprule
\textbf{set} & \textbf{k} & \textbf{$\chi^2$} & \textbf{$p$} & \textbf{Kendall $W$} & \textbf{rejects} \\
\midrule
headline (Stage A+B) & 7 & 4.52 & 0.607 & 0.027 & \textbf{no} \\
segformer\_family & 3 & 0.93 & 0.629 & 0.017 & \textbf{no} \\
mobile\_field & 5 & 16.26 & 0.0027 & 0.145 & yes \\
kd\_arms (Stage C) & 10 & 20.63 & 0.0144 & 0.082 & yes \\
\bottomrule
\end{tabular}
\end{center}

\begin{keybox}{The most defensible sentence in the study}
\textbf{The seven headline models are collectively indistinguishable at subject
level} ($p=0.607$, Kendall $W=0.027$). This is \S1's annotation-ceiling finding
arriving through a second, independent route --- and a stronger form of it,
because it is a formal test rather than a comparison of magnitudes. $W=0.027$
means subjects essentially do not agree on how to rank these models.
\end{keybox}

The Stage C arms do reject, but read that carefully: those ten arms include
several that are plainly weaker (\texttt{x\_angular} 0.7418,
\texttt{p3\_adaptive\_hard} 0.7414 against \texttt{p3\_adaptive} 0.7748), so the
omnibus is detecting the bad arms, not vindicating the good ones. It licenses
pairwise comparisons \emph{inside} Stage C and says nothing about the headline
field.

\subsection{The contrast family}

Every endpoint comes from the \emph{same} resampled subject lists (10{,}000
draws), so ``Dice up, misses up'' is a statement about the same resampled worlds
rather than three unrelated bootstraps.

\begin{center}
\small\setlength{\tabcolsep}{3.5pt}
\begin{tabular}{p{4.05cm}p{3.05cm}lnp{2.85cm}l}
\toprule
\textbf{A} & \textbf{B} & \textbf{kind} & \textbf{$\Delta$} & \textbf{95\,\% CI} & \textbf{verdict} \\
\midrule
lraspp\_mnv3\_distilled & lraspp\_mnv3 & conf & $+0.0218$ & $[+0.0081,+0.0372]$ & \win \\
segformer\_b0\_distilled & segformer\_b0\_direct & conf & $+0.0017$ & $[-0.0088,+0.0135]$ & \noninf \\
segformer\_b0\_distilled & segformer\_b2\_teacher & conf & $-0.0012$ & $[-0.0197,+0.0179]$ & \inc \\
segformer\_b2\_teacher & unet\_r50 & expl & $+0.0122$ & $[-0.0263,+0.0515]$ & \inc \\
unet\_r50 & deeplabv3plus\_r50 & expl & $-0.0014$ & $[-0.0184,+0.0144]$ & \inc \\
segformer\_b0\_direct & yolo\_sem\_direct & expl & $+0.0642$ & $[+0.0128,+0.1200]$ & \win \\
lraspp\_mnv3 & topformer\_tiny & expl & $+0.0063$ & $[-0.0439,+0.0552]$ & \inc \\
lraspp\_mnv3 & fastscnn & expl & $+0.0928$ & $[+0.0209,+0.1572]$ & \win \\
segformer\_b0\_direct & lraspp\_mnv3 & expl & $+0.0682$ & $[+0.0244,+0.1175]$ & \win \\
lraspp\_mnv3\_distilled & deeplabv3plus\_r50 & expl & $-0.0384$ & $[-0.0912,+0.0027]$ & \inc \\
\bottomrule
\end{tabular}
\end{center}

Two contrasts were \textbf{skipped and named} because
\texttt{fastscnn\_distilled} has no current-lineage result. Holm therefore ran
over three confirmatory contrasts rather than four; at $\times4$ the LR-ASPP
adjusted $p$ would be 0.0056, so that verdict does not depend on it.

\subsection{INCONCLUSIVE is not NON-INFERIOR}
\label{sec:inconclusive}

Stage C's original rule was WIN / INFERIOR / \emph{anything else}
$\rightarrow$ NON-INFERIOR. That silently classes an interval running from $-0.05$
to $+0.02$ as equivalence, which reports absence of evidence as evidence of
absence. The corrected rule is four-way:

\begin{center}
\begin{tabular}{ll}
\toprule
\win & $\mathrm{lo} > 0$ \\
\infr & $\mathrm{hi} < -\text{margin}$ \\
\noninf & $\mathrm{lo} \ge -\text{margin}$ \\
\inc & interval extends below $-$margin but not entirely \\
\bottomrule
\end{tabular}
\end{center}

\textbf{This changes real verdicts.} \texttt{segformer\_b0\_distilled} vs
\texttt{segformer\_b2\_teacher} is \inc\ (CI $[-0.020,+0.018]$), not the
equivalence the old rule would have called it. \textbf{Stage C's ten arms should be
re-scored under this rule before their NON-INFERIOR verdicts are quoted again.}

\subsection{Complete misses, as counts}

At 0--13 misses of 185 a bootstrapped difference of two rates near the boundary is
unstable. The $2\times2$ is the evidence itself.

\begin{center}
\small\setlength{\tabcolsep}{4pt}
\begin{tabular}{p{3.85cm}p{3.15cm}nnnl}
\toprule
\textbf{A} & \textbf{B} & \textbf{A only} & \textbf{B only} & \textbf{both} & \textbf{containment} \\
\midrule
segformer\_b2\_teacher & unet\_r50 & 0 & 7 & 0 & \textbf{b2 within unet} \\
segformer\_b0\_direct & yolo\_sem\_direct & 0 & 11 & 1 & \textbf{b0 within yolo} \\
segformer\_b0\_distilled & segformer\_b0\_direct & 0 & 1 & 0 & distilled within direct \\
unet\_r50 & deeplabv3plus\_r50 & 2 & 0 & 5 & deeplab within unet \\
lraspp\_mnv3 & fastscnn & 2 & 13 & 0 & none \\
lraspp\_mnv3\_distilled & lraspp\_mnv3 & 2 & 1 & 1 & none \\
\bottomrule
\end{tabular}
\end{center}

\begin{keybox}{Where the headline field actually separates}
The omnibus says the seven headline models are indistinguishable on Dice. On
complete misses they separate by \textbf{strict containment}: no model in the
SegFormer line ever misses a bruise its comparator finds. That is a far stronger
statement than any Dice contrast in the study, and it comes out of counts rather
than intervals.
\end{keybox}

% ============================================================================
\section{The complete-miss definition, and a resolved measurement bug}

For months two artefacts disagreed on Stage E miss counts and the discrepancy was
recorded as unexplained. It was resolved on 30 July 2026, and \emph{neither
artefact was corrupt}.

\begin{center}
\begin{tabular}{Lnnn}
\toprule
\textbf{model} & \textbf{per-seed CSV} & \textbf{pred\_pos = 0} & \textbf{dice = 0} \\
\midrule
topformer\_tiny & 0 & \textbf{0} & 1 \\
ppmobileseg\_tiny & 3 & \textbf{3} & 4 \\
fastscnn & 8 & \textbf{8} & 13 \\
lraspp\_mobilenetv3 & 2 & \textbf{2} & 2 \\
\bottomrule
\end{tabular}
\end{center}

\texttt{efficient\_test\_per\_seed.csv} counts images where the model
\textbf{predicted nothing at all}. \texttt{report.normalize} counts images where
the prediction has \textbf{no overlap with the ground truth}. The first is a
strict subset of the second: it excludes the case where a model fires confidently
on the wrong region --- which is still a complete miss to a clinician, and is
exactly the failure the metric exists to catch. LR-ASPP's row matched only because
it happens to have no fire-in-the-wrong-place cases.

\begin{keybox}{Publish the \texttt{dice = 0} column}
Fast-SCNN's true complete-miss rate is 13/185 = \textbf{7.0\,\%}, not the 4.3\,\%
the per-seed table implies. Every Stage E miss figure sourced from the per-seed
table understates by the same mechanism. Neither file needs ``fixing'' --- they
measure different things and only one of them is the clinical metric.
\end{keybox}

% ============================================================================
\section{The final ranking}

Sorted by mean Dice at the val-selected seed. \texttt{(C)} = Stage C arm, single
seed 42. \texttt{fastscnn\_distilled} is excluded --- no current-lineage result.
Miss counts are $\mathrm{dice}=0$.

\begin{center}
\small
\begin{longtable}{rp{5.6cm}nnnn}
\toprule
\textbf{\#} & \textbf{model} & \textbf{mean} & \textbf{median} & \textbf{miss} & \textbf{miss \%} \\
\midrule
\endhead
1 & p3\_adaptive (C) & 0.7748 & 0.8172 & 2 & 1.08 \\
2 & p2\_cwd\_b5\_to\_b0 (C) & 0.7736 & 0.8230 & 4 & 2.16 \\
3 & \textbf{segformer\_b5\_teacher} & 0.7727 & 0.8273 & \textbf{0} & \textbf{0.00} \\
4 & \textbf{segformer\_b2\_teacher} & 0.7692 & 0.8192 & \textbf{0} & \textbf{0.00} \\
5 & expA\_b5\_to\_b0\_response (C) & 0.7683 & 0.8061 & 1 & 0.54 \\
6 & \textbf{segformer\_b0\_distilled} & 0.7680 & 0.8167 & \textbf{0} & \textbf{0.00} \\
7 & segformer\_b5\_\_seed2026 (C) & 0.7676 & 0.8298 & 1 & 0.54 \\
8 & p3\_adaptive\_boundary (C) & 0.7668 & 0.8221 & 1 & 0.54 \\
9 & segformer\_b0\_direct & 0.7663 & 0.8129 & 1 & 0.54 \\
10 & p2\_bpkd\_b5\_to\_b0 (C) & 0.7651 & 0.8091 & 1 & 0.54 \\
11 & p2\_ensemble\_uniform (C) & 0.7615 & 0.8139 & 2 & 1.08 \\
12 & p3\_adaptive\_group (C) & 0.7586 & 0.8047 & 2 & 1.08 \\
13 & deeplabv3plus\_r50 & 0.7584 & 0.8183 & 5 & 2.70 \\
14 & unet\_r50 & 0.7570 & \textbf{0.8329} & 7 & 3.78 \\
15 & segformer\_b5\_\_seed42 (C) & 0.7527 & 0.8124 & 3 & 1.62 \\
16 & p3\_adaptive\_full (C) & 0.7515 & 0.7984 & 3 & 1.62 \\
17 & x\_angular\_b5\_to\_b0 (C) & 0.7418 & 0.8043 & 2 & 1.08 \\
18 & p3\_adaptive\_hard (C) & 0.7414 & 0.8023 & 1 & 0.54 \\
19 & yolo\_sem\_distilled & 0.7261 & 0.8012 & 5 & 2.70 \\
20 & lraspp\_mobilenetv3\_distilled & 0.7200 & 0.7836 & 3 & 1.62 \\
21 & yolo\_sem\_direct & 0.7021 & 0.8061 & 12 & 6.49 \\
22 & lraspp\_mobilenetv3 & 0.6982 & 0.7607 & 2 & 1.08 \\
23 & topformer\_tiny & 0.6918 & 0.7418 & \textbf{1} & \textbf{0.54} \\
24 & ppmobileseg\_tiny & 0.6568 & 0.7158 & 4 & 2.16 \\
25 & fastscnn & 0.6053 & 0.6866 & 13 & 7.03 \\
\bottomrule
\end{longtable}
\end{center}

\subsection{Read it as three tiers, not twenty-five places}

\textbf{Most of this ordering is not real.} Ranks 1--12 are separated by 0.013
Dice in total --- less than a fifth of the human-vs-human disagreement --- and the
omnibus does not reject for the headline set.

\begin{sloppypar}
\textbf{Tier 1 --- zero complete misses (3 models):}
\texttt{segformer\_b5\_teacher}, \texttt{segformer\_b2\_teacher},
\texttt{segformer\_b0\_distilled}. This is the only cut in the table that is a
real distinction rather than a Dice ordering.

\textbf{Tier 2 --- usable, 1--2 misses:} \texttt{segformer\_b0\_direct}, the
better Stage C arms, \texttt{topformer\_tiny}.

\textbf{Tier 3 --- miss too often for a find-the-bruise task:}
\texttt{unet\_r50} (7), \texttt{yolo\_sem\_direct} (12), \texttt{fastscnn} (13).
\end{sloppypar}

\subsection{The recommendation}

\begin{keybox}{What to ship}
\textbf{Research or internal use: \texttt{segformer\_b0\_distilled}.} Zero
complete misses, 0.7680 mean Dice, 3.71\,M parameters, 16.6\,ms. Statistically
indistinguishable from its own 27\,M teacher, so nothing gained by going bigger is
real. \textbf{B2 $\rightarrow$ B0 is $7.4\times$ fewer parameters for $-0.005$
Dice} --- the strongest clean compression result in the study.

\textbf{Product / commercial: \texttt{deeplabv3plus\_r50}} if 26\,M parameters is
affordable (5 misses), otherwise \textbf{\texttt{lraspp\_mobilenetv3\_distilled}}
(3 misses, $-0.048$ Dice). Every Tier-1 model is SegFormer, and MiT weights are
NVIDIA non-commercial.

\textbf{If complete-miss rate is the only thing that matters and the budget is
mobile: \texttt{topformer\_tiny}} --- 1 miss at 1.37\,M parameters, the best
miss rate in the mobile field despite ranking below both LR-ASPP variants on Dice.
\end{keybox}

\subsection{Is smaller-but-equal possible? Yes --- but not through parameter count}

The B2 $\rightarrow$ B0 step is $7.4\times$ compression for $-0.005$ Dice, so the
answer is plainly yes. What the mobile results show is that \textbf{parameter
count is not the binding constraint}: LR-ASPP is 3.22\,M --- 13\,\% \emph{fewer}
parameters than B0 --- and is 0.05 Dice behind. YOLO at 1.63\,M holds a median
within 0.007 of B0 while missing 12 images.

\begin{center}
\begin{tabular}{Lnnn}
\toprule
\textbf{model} & \textbf{params} & \textbf{median Dice} & \textbf{miss \%} \\
\midrule
segformer\_b0\_direct & 3.71\,M & 0.8129 & 0.54 \\
yolo\_sem\_direct & \textbf{1.63\,M} & 0.8061 & 6.49 \\
lraspp\_mobilenetv3 & 3.22\,M & 0.7607 & 1.08 \\
\bottomrule
\end{tabular}
\end{center}

\textbf{Median accuracy survives compression down to 1.6\,M almost untouched;
what breaks is the tail.} ``Preserving accuracy'' splits in two, and only one half
is cheap to keep. The mobile models buy \emph{latency}, not size.

Caveat: the size question has not actually been tested. The smallest SegFormer
here is B0 --- there is no half-width MiT, no quantisation, no pruning, no
lower-resolution variant. What is tested is that five off-the-shelf mobile
architectures are worse than B0.

% ============================================================================
\section{Traps hit, and the guards against them}

Each of these was a real mistake with a real cost. Each now has an automated
check.

\begin{enumerate}[leftmargin=*,itemsep=3pt]
\item \textbf{A same-named checkpoint set with the wrong lineage.} A directory
with all fifteen Stage A run\_ids but $\alpha=0.5$, batch 8, and YOLO trained in
the custom loop. \emph{Guard:} the build asserts on $\alpha=0.6$, per-model batch
and the native Ultralytics weight path. Names are not evidence; configs are.

\item \textbf{A baseline set from the wrong split} (693/138 at seed 42, not the
canonical 697/134 three-seed set). \emph{Guard:} the build reconciles each run's
own per-image CSV against the canonical per-seed CSV and refuses mismatches.

\item \textbf{The best seed is not the same for every model} --- 0 for the three
SegFormers, 2 for \texttt{yolo\_sem\_direct}. Pairing weights with another seed's
results shows per-image disagreements up to 0.49 Dice and looks exactly like
broken inference. \emph{Guard:} \texttt{report.best\_seeds()} reads the mapping
off the selection step's own filenames.

\item \textbf{The same model under two names.} \texttt{segformer\_b5\_teacher} is
seed 123 promoted; listing both showed one result twice. \emph{Guard:}
\texttt{load\_stage\_c} fingerprints the per-image Dice vector and collapses exact
duplicates.

\item \textbf{A silent ImageNet download.} \texttt{smp} fetches ResNet-50 the
moment a model is constructed, even when a checkpoint is about to overwrite every
weight. \emph{Guard:} \texttt{encoder\_weights=None} on the load path.

\item \textbf{An activation substitution that changes nothing visible} ---
\texttt{nn.Hardsigmoid} for \texttt{HSigmoid(slope=0.2)}. Every state-dict key
still matches while every SE gate computes something different. \emph{Guard:}
the shim honours slope/offset, and \texttt{verify\_checkpoint\_match} proves the
load.

\item \textbf{Autograd graph retained during evaluation.} Invisible on a 40\,GB
A100; on CPU, B2's decode head dies allocating 1.2\,GB in one \texttt{cat}.
\emph{Guard:} \texttt{torch.inference\_mode()} at the call site.

\item \textbf{Bootstrap by DataFrame concatenation} --- 2000 \texttt{pd.concat}
calls per model $\times$ 16 models = 32k allocations, enough to kill a Colab
kernel. \emph{Guard:} resample index arrays instead.

\item \textbf{Eager dependency installs on an offline node.} \emph{Guard:}
optional installs gated on whether a model will actually be \emph{built}.

\item \textbf{Demanding CUDA for work that will not happen.} \emph{Guard:} filter
to trainable kinds first, then require CUDA.

\item \textbf{Writing a cell instead of following the tested pattern.}
\emph{Lesson:} when a tested cell for the same job exists, copy its call
sequence.

\item \textbf{Reading a 0-of-185 result as a zero rate.} LR-ASPP's seed 0 missed
nothing and this was reported as a zero-miss model; seeds 1 and 2 missed 1 and 2.
\emph{Lesson:} a zero count on $n=185$ has a one-sided 95\,\% bound of
$\approx1.6\,\%$, which spans the entire Stage E field.

\item \textbf{Documented results read as shipped artefacts.} A results table in a
document does not mean the weights exist. \emph{Guard:} the registry's tier table
is the authority, not a results section.

\item \textbf{Two definitions of the same metric.} Resolved above --- empty
prediction vs zero overlap.

\item \textbf{Calling an underpowered interval ``equivalence.''} Resolved by the
four-way verdict rule.
\end{enumerate}

% ============================================================================
\section{Known gaps and limitations}

\subsection{Genuine gaps}
\begin{enumerate}[leftmargin=*,itemsep=2pt]
\item \textbf{nnU-Net} --- never run on the canonical 697/134 split. No weights,
no results. Needs \texttt{nnunetv2} and $\sim$8 GPU-hours.
\item \textbf{\texttt{x\_dkd\_b5\_to\_b0}} --- a KD arm configured but never
executed; the directory holds only \texttt{run\_config.json}.
\item \textbf{\texttt{fastscnn\_distilled}} --- no current-lineage result at any
seed. $\sim$2 GPU-hours.
\end{enumerate}

All three appear in \texttt{gaps.csv}. None is back-filled.

\subsection{Limitations that are not gaps}
\begin{enumerate}[leftmargin=*,itemsep=2pt,start=4]
\item \textbf{Stage E/F ship no weights in the bundle.} The runs live on the ORC
cluster work directory; the registry correctly reports them MISSING from a local
bundle, which is easy to misread as needless retraining.
\item \textbf{B5 seeds 42 and 2026 have results but no weights} --- only the
val-selected seed was promoted. Their numbers are exact; they cannot be
re-inferred.
\item \textbf{\texttt{yolo\_sem\_distilled\_\_seed2} stopped at 31 epochs} where
its siblings ran 49--84, and its \texttt{best.pt} was never stripped of optimizer
state. It loads and scores correctly; it is why that model's seed spread is wide
(std 0.065). Preserved as-is rather than silently re-run.
\item \textbf{Speed tables span two devices} (full A100 vs MIG 3g.40gb). Not
comparable.
\item \textbf{TopFormer's weights need a manual download.} Without them it trains
from scratch and is labelled as such.
\item \textbf{28 test subjects is a small denominator.} Most fairness comparisons
are not significant and must not be reported as if they were.
\item \textbf{Stage C's ten arms are single-seed (42) only} --- no seed spread
exists for any of them.
\item \textbf{Stage F fairness was never computed for the LR-ASPP arm}, despite
being its pre-registered primary endpoint.
\item \textbf{transformers version pinning.} Recent releases renamed SegFormer's
internals; Stage A checkpoints will not load across that refactor.
\end{enumerate}

% ============================================================================
\section{Next steps}
\label{sec:next}

In priority order. Only items 1 and 5 cost GPU time.

\subsection*{1. Re-run \texttt{fastscnn\_distilled} at three seeds ($\sim$2 GPU-h)}
\textbf{This is the single highest-value action available.} The two-arm
student-initialisation finding --- that cross-architecture KD hurts a scratch
student and helps a pretrained one --- is the most interesting claim Stage F can
make, and it is currently a cross-lineage comparison. One arm is current, one is
from a superseded sweep. Two GPU-hours converts an anecdote into a controlled
two-point design, and restores Holm correction to its pre-specified family of
four.

\subsection*{2. Compute Stage F fairness (CPU, minutes)}
The LR-ASPP arm was pre-registered as a miss-rate \emph{and fairness} claim.
Fairness has never been run on it. Given that the Fast-SCNN arm widened its ITA
gap from 0.136 to 0.199 at $p=0.0017$, whether the LR-ASPP arm tightened or
widened its own gap is a genuinely open question with a real chance of changing
the story. Run \texttt{report.fairness\_by\_group} and
\texttt{report.fairness\_gap} on all three seeds.

\subsection*{3. Re-score Stage C's ten arms under the four-way verdict rule (CPU)}
Every arm currently carries a NON-INFERIOR verdict assigned by a rule that could
not distinguish equivalence from an underpowered interval. At least one
reference contrast flips to INCONCLUSIVE under the corrected rule. Re-running
Stage G with the arms in the contrast family settles it in minutes.

\subsection*{4. Run the teacher-miss containment check for the LR-ASPP arm (CPU)}
The arm's misses went from 3 to 5 across seeds while its teacher misses more than
it does. Whether the student's new misses fall on images the teacher also misses
is the difference between ``KD imported a teacher weakness'' and ``random
variation.'' The paired $2\times2$ against \texttt{deeplabv3plus\_r50} answers it
directly and is already implemented.

\subsection*{5. nnU-Net on the canonical split ($\sim$8 GPU-h)}
The one named gap in the baseline set, and the standard medical-segmentation
baseline a reviewer will ask about. It is the only remaining training run whose
absence is a reasonable objection to the study rather than a scoping decision.

\subsection*{6. Size-conditioned fairness, properly (CPU)}
Bruise size confounds every fairness claim in the study, including the
counter-intuitive Light-skin YOLO result. The stratified and regression analyses
exist but are not yet the primary presentation of fairness. They should be.

\subsection*{7. Re-time Stage A and Stage E on one device (GPU-minutes)}
Two speed tables on two devices cannot be compared, and the accuracy-vs-latency
trade --- the actual deployment question --- cannot be stated until they can.

\subsection*{8. What NOT to do}
\textbf{Do not add a sixth mobile architecture.} Every model is at the annotation
ceiling; another one produces another row inside the noise floor and answers
nothing. The study's remaining value is in making the existing 52 runs airtight,
not in adding a 53rd.

\begin{keybox}{The three contributions that need no new training}
\begin{enumerate}[leftmargin=*,itemsep=1pt]
\item \textbf{The annotation ceiling}, now backed by a formal omnibus test that
does not reject across the headline field.
\item \textbf{B2 $\rightarrow$ B0: $7.4\times$ compression for $-0.005$ Dice},
with zero complete misses preserved.
\item \textbf{Complete-miss containment} as the metric that actually separates
these models, where Dice does not.
\end{enumerate}
Item 1 of \S\ref{sec:next} adds a fourth: KD's direction depends on student
initialisation.
\end{keybox}

% ============================================================================
\appendix
\section{Where every number came from}

\begin{center}
\small
\begin{tabular}{p{4.9cm}>{\raggedright\arraybackslash}p{9.3cm}}
\toprule
\textbf{Table / claim} & \textbf{Source} \\
\midrule
Annotation ceiling & \texttt{interlabeler\_agreement\_640.csv}; \texttt{results/analysis\_native/annotation\_ceiling.csv} \\
Stage A 3-seed results & \texttt{results/final/FINAL\_RESULTS.csv}, per-seed CSVs \\
Stage A speed & \texttt{results/final/benchmark\_640.csv} \\
Stage A fairness & \texttt{results/analysis\_native/fairness\_stats\_bestseed.csv} \\
Stage B results & \texttt{results/baselines/BASELINES\_RESULTS.csv} \\
Stage C teacher, oracle, alpha, arms & \texttt{checkpoints/distill/}, \texttt{results/distill/}, \texttt{stage\_c\_vs\_reference.csv} \\
Stage E results & \texttt{\_work/new\_outputs/headline\_stage\_e.csv}, \texttt{efficient\_test\_per\_seed.csv} \\
Stage E checkpoint match & \texttt{efficient\_checkpoint\_match.csv} \\
Stage F LR-ASPP per-seed & ORC \texttt{WORK\_DIR/runs/} \texttt{lraspp\_mobilenetv3\_distilled\_\_seed\{0,1,2\}} \\
Stage G omnibus & \texttt{significance\_omnibus.csv} \\
Stage G contrast family & \texttt{significance\_contrast\_family.csv} \\
Stage G discordance & \texttt{significance\_miss\_discordance.csv} \\
Stage G seed consistency & \texttt{significance\_seed\_consistency.csv} \\
Final ranking & recomputed from all \texttt{per\_image\_*.csv} in \texttt{\_work/new\_outputs/} \\
Miss-definition table & recomputed: \texttt{(dice==0)} vs \texttt{(pred\_positive\_pixels==0)} \\
Gaps & \texttt{gaps.csv} \\
\bottomrule
\end{tabular}
\end{center}

\vspace{8pt}
\begin{sloppypar}
\noindent\textbf{Reproducing everything:} the analysis is a function of per-image
Dice, not of model weights. \texttt{bruise\_unified.ipynb} with defaults
(\texttt{ALLOW\_TRAINING=False}, \texttt{RECOMPUTE\_FROM\_WEIGHTS=False})
regenerates every table and figure in this document in about two minutes on a
laptop with no GPU. Full method documentation is in
\texttt{PROJECT\_HANDBOOK.md}.
\end{sloppypar}

\end{document}
